\section{Related Work}

\subsection{CoT Faithfulness and Monitoring}

Recent work has revealed significant concerns about the faithfulness of chain-of-thought (CoT) reasoning in language models. \citet{cot-wild-2025} demonstrate that models can produce logically contradictory arguments, generating superficially coherent but fundamentally inconsistent reasoning through what they term ``Implicit Post-Hoc Rationalization.'' Similarly, \citet{cot-monitors-2025} investigate CoT monitoring as an AI safety mechanism, distinguishing between ``CoT-as-rationalization'' and ``CoT-as-computation'' and showing that models can learn to obscure their true reasoning processes. Our work complements these findings by providing a systematic method to detect when models make decisions based on unverbalized concepts, helping identify cases where explanations may not faithfully represent the underlying decision-making process.

Several works have developed methods to evaluate whether model explanations accurately reflect decision-making processes. \citet{faithfulness-tests-2023} propose counterfactual input editors and reconstruction methods to test the faithfulness of natural language explanations. \citet{probabilities-matter-2024} introduce the Correlational Counterfactual Test (CCT), which considers the total shift in predicted label distributions rather than just binary outcomes. \citet{causal-lens-2025} develop the Causal Diagnosticity framework for evaluating faithfulness metrics, finding that continuous metrics are generally more diagnostic than binary ones. Our statistical testing approach aligns with these findings and provides a complementary method by identifying concepts that show significant performance differences but are not cited as justification in model outputs.

\citet{llm-decision-boundaries-2025} demonstrate that LLMs cannot reliably generate minimal counterfactual explanations, producing either overly verbose or insufficiently modified inputs. This unreliability of self-generated explanations motivates external testing approaches like ours, which do not rely on the model's own introspection capabilities.

\subsection{Bias Detection via Counterfactuals}

Hidden biases in language models have been documented across several applications. \citet{llm-fairness-2025} reveal that models exhibit significant demographic biases in hiring applications even with subtle contextual cues, showing that models can infer sensitive attributes from indirect information. \citet{prefix-bias-2025} investigate how minor variations in query prefixes can systematically shift model preferences across race and gender dimensions. \citet{implicit-bias-llm-2024} adapt the Implicit Association Test from psychology to LLMs, revealing pervasive implicit stereotype biases across race, gender, religion, and health categories in eight value-aligned models that appear unbiased on standard benchmarks. These works rely on manually curated datasets or predefined categories targeting specific dimensions. Our approach is fully automatic: we generate both concept hypotheses and test variations on the fly, extending beyond predefined categories to systematically discover a broader range of bias-inducing concepts. Furthermore, we explicitly check whether detected biases appear in the model's chain-of-thought reasoning, targeting unverbalized biases.